\documentclass[12pt,a4paper]{article}

% =================================================
% Page layout (A4, scientific margins)
% =================================================
\usepackage[a4paper,
            left=25mm,
            right=25mm,
            top=25mm,
            bottom=25mm]{geometry}

% =================================================
% Fonts and justification
% =================================================
\usepackage[T1]{fontenc}
\usepackage{lmodern}
\usepackage{setspace}
\usepackage{ragged2e}
\usepackage{microtype}

% =================================================
% Code listings
% =================================================
\usepackage{listings}
\usepackage{xcolor}

\lstset{
  basicstyle=\ttfamily\small,
  frame=single,
  breaklines=true,
  columns=fullflexible,
  showstringspaces=false
}

% =================================================
% Hyperlinks
% =================================================
\usepackage[hidelinks]{hyperref}

% =================================================
% Document begins
% =================================================
\begin{document}

\begin{center}
    {\LARGE \textbf{Jobsubmit User Manual}} \\[0.4cm]
    {\large Version 1.0.8} \\[0.2cm]
    {\large GNU General Public License (GPL)} \\[1cm]
    \textbf{Author} \\[0.2cm]
    Rajesh Prashanth Anandavadivel
\end{center}

\vspace{1cm}
\onehalfspacing
\justifying

% =================================================
\section{Introduction}

Jobsubmit is a lightweight, platform-independent job execution utility designed for scientific computing, batch processing, and reproducible workflows. It executes commands from a list file sequentially while clearly reporting job start time, completion status, execution duration, process ID, and exit code.

The design philosophy emphasizes simplicity, transparency, and deterministic behavior, making Jobsubmit suitable for local machines, login nodes, virtual machines, continuous integration pipelines, and development environments.

Jobsubmit is not intended to replace full-featured workload managers or schedulers. Instead, it provides a minimal and predictable execution layer for controlled, sequential workflows.

% =================================================
\section{Key Features}

Jobsubmit provides the following core features:

\begin{itemize}
    \item Sequential execution of commands from a text file
    \item Explicit reporting of:
    \begin{itemize}
        \item Job start date and time
        \item Job completion or failure date and time
        \item Execution duration in HH:MM:SS format
        \item Process ID (PID)
        \item Exit code
    \end{itemize}
    \item Separate capture of standard output and standard error
    \item Clean, scientific-style console output
    \item Cross-platform compatibility (Linux, macOS, Windows)
    \item No external Python dependencies
    \item Simple installation and clean project layout
\end{itemize}

% =================================================
\section{Project Structure}

The Jobsubmit distribution follows a simple and transparent directory layout:

\begin{lstlisting}
jobsubmit/
│   .gitlab-ci.yml
│   .travis.yml
│   CHANGELOG
│   LICENSE
│   README.md
│
├── bin/
│   └── jobsubmit
│
├── docs/
│
├── logs/
│
├── tests/
│   └── cmd.lst
│
└── tools/
    └── install.sh
\end{lstlisting}

% =================================================
\section{System Requirements}

\subsection{Software Requirements}
\begin{itemize}
    \item Python version 3.7 or newer
    \item Bash shell (required for the installer on Linux and macOS)
\end{itemize}

No third-party Python modules are required.

\subsection{Hardware Requirements}
Jobsubmit has no specific hardware requirements and can be executed on desktops, laptops, virtual machines, and CI runners.

% =================================================
\section{Installation}

\subsection{Interactive Installation (Linux / macOS)}

From the project root directory, execute:

\begin{lstlisting}
bash tools/install.sh
\end{lstlisting}

The installer performs the following actions:

\begin{itemize}
    \item Prompts for an installation directory
    \item Copies the Jobsubmit executable
    \item Sets executable permissions
    \item Optionally adds the installation directory to the user's PATH
    \item Creates output directories for standard output and error
\end{itemize}

After installation, reload the shell configuration:

\begin{lstlisting}
source ~/.bashrc
\end{lstlisting}

\subsection{Manual Installation}

\begin{lstlisting}
mkdir -p ~/jobsubmit
cp bin/jobsubmit ~/jobsubmit/
chmod +x ~/jobsubmit/jobsubmit
export PATH=$PATH:~/jobsubmit
\end{lstlisting}

To make the PATH modification permanent, add the export command to \texttt{\textasciitilde/.bashrc} or \texttt{\textasciitilde/.profile}.

\subsection{Windows Installation}

On Windows systems, Jobsubmit can be executed directly using Python:

\begin{lstlisting}
python bin\jobsubmit cmd.lst all
\end{lstlisting}

Executable binaries generated using tools such as Nuitka or PyInstaller may also be used if available.

% =================================================
\section{Usage}

\subsection{Command Syntax}

\begin{lstlisting}
jobsubmit <command_file> <line_selection>
\end{lstlisting}

\subsection{Examples}

\begin{lstlisting}
jobsubmit cmd.lst all
jobsubmit cmd.lst 3
jobsubmit cmd.lst 1-5
jobsubmit cmd.lst 1,3,7
\end{lstlisting}

% =================================================
\section{Command File Format}

The command file (commonly named \texttt{cmd.lst}) defines the workflows to be executed.

\begin{itemize}
    \item One command per line
    \item Blank lines are ignored
    \item Commands are executed sequentially
\end{itemize}

\subsection{Example}

\begin{lstlisting}
echo "Job 1"
sleep 2
invalid_command
echo "Job 4"
\end{lstlisting}

% =================================================
\section{Output Files}

For each executed command, Jobsubmit generates:

\begin{itemize}
    \item \texttt{stdout/<cmdfile>\_line\_<N>.out}
    \item \texttt{stderr/<cmdfile>\_line\_<N>.err}
\end{itemize}

These files enable post-run inspection, debugging, and archival.

% =================================================
\section{Exit Codes}

\begin{itemize}
    \item Exit code \texttt{0}: Command executed successfully
    \item Non-zero exit code: Command execution failed
\end{itemize}

Jobsubmit itself exits with a non-zero status only in the event of fatal errors such as invalid input or missing files.

% =================================================
\section{Continuous Integration}

Jobsubmit includes configuration files for:

\begin{itemize}
    \item GitLab CI (\texttt{.gitlab-ci.yml})
    \item Travis CI (\texttt{.travis.yml})
\end{itemize}

These configurations ensure consistent behavior across supported platforms.

% =================================================
\section{License}

Jobsubmit is distributed under the \textbf{GNU General Public License (GPL)}.
Users are free to use, modify, and redistribute the software under the terms of the GPL.

See the \texttt{LICENSE} file for the complete license text.

% =================================================
\section{Author}

\textbf{Rajesh Prashanth Anandavadivel}

Originally developed as a practical utility for managing reproducible scientific and computational workflows.

% =================================================
\section{Changelog}

All notable changes, version updates, and fixes are documented in the \texttt{CHANGELOG} file.

% =================================================
\section{Philosophy}

Jobsubmit is intentionally:

\begin{itemize}
    \item Minimal
    \item Predictable
    \item Transparent
    \item Script-friendly
\end{itemize}

It is not intended to replace full-fledged schedulers such as SLURM or PBS, but to provide a reliable execution tool for controlled, sequential workflows.

% =================================================
\section{Contributing}

Bug reports, improvements, and suggestions are welcome.
Contributions should remain focused, readable, and aligned with the scientific intent of the project.

\vfill
\begin{center}
\textit{End of User Manual}
\end{center}

\end{document}
